\label{sec:intro}

The introduction should give the reader the neccessary background. For
example, start with a bird's perspective and zoom in on the smaller
details investigated in the present project.

Citations supporting your statements look like this:
\cite{artref,procref}, that is, a cited article and a conference
proceeding. Single book citations \cite{bookref} or master's theses
\cite{thesisref} may also be used. Frequently, so-called
\emph{preprints} are cited before they even have been published
\cite{preprintref}. It goes without saying that one might also want to
cite web-pages, like a specific URL \cite{urlref} or, perhaps more
common, a Wikipedia-page \cite{wikiref}.

\DO{Remember that the citations are part of your sentence and should
  be included in the sentence as such, i.e., come before a `.' at the
  end of the sentence, are separated by commas where appropriate,
  etc.}  \DO{Look at the Bibliography file and/or check any Bib\TeX\
  warnings and make sure to fill out all the appropriate fields for
  the type of reference you use.}

\DONT{Citations should be used to support and reference what you have
  written. Readers should understand the main point you have taken
  from the content of your citation without having to read it
  themselves.}

\DO{ It is common to end the introduction with an overview of the
  structure of the report.} In this example, theory is summarized in
\secref{theory}, and the required methods in \secref{methods}. The
results are to be found in \secref{results} and a discussion around
the central themes of the project is finally offered in
\secref{conclusions}. Note the use of \verb!\secref{<section label>}!
to reference sections.